\section{Proof ideas: obstacles and the estimates that overcome them}
\label{sec:proof-ideas}

The two halves of the theorem require different information. To prove escape,
one must show that a small disk cannot compete with the global local-time
maximum. To prove critical recurrence, one must bring the entire set of
maximizers into a small disk despite all local time accumulated in the past.
We explain the quantitative mechanisms before developing their proofs.

\subsection{Why a site-by-site bound misses the escape threshold}

The local time at a site, after the first hit and before leaving a disk of
radius $e^J$, has a geometric tail with mean $gJ+O(1)$. At level $u_{J,c}$
this gives
\[
 P_x(L_{\tau_{e^J}}(x)\ge u_{J,c})
 \le C\exp\{-u_{J,c}/(gJ)+O(1)\}
 =Ce^{-2J}J^c
\]
for sites in a fixed inner fraction of the disk. A union bound over
$e^{2J}J^{-2\beta}$ sites therefore gives only $CJ^{c-2\beta}$.
Even incorporating the hitting probability does not give the needed
estimate. In a logarithmic shell at distance about $e^{J-k}$, that
probability is at most $C(k+1)/J$, and the number of sites is
$O(e^{2J-2k})$. Summing over $k\ge\beta\log J+O(1)$ yields
\begin{equation}\label{eq:ideas-naive-upper}
 C J^{c-1}\sum_{k\ge\beta\log J+O(1)}(k+1)e^{-2k}
 \le J^{-2\beta+c-1+o(1)}.
\end{equation}
At the relevant level $c\downarrow2$, summability of this bound over
exponential time blocks would require $\beta>1$. The desired conclusion
requires the extra factor $J^{-1}$ in
\eqref{eq:intro-localized-upper}.

The improvement comes from recording how local time is built across nested
annuli. Let $T_j$ count inward traversals at logarithmic level $j$.
For the ideal concentric Brownian model, descendants of completed source
excursions form a critical Galton--Watson process with geometric offspring.
The square-root coordinate $\sqrt{2T_j}$ has an approximately linear
profile on an extreme traversal event. We restrict this profile below
\[
 B_j=\rho_Jj+z+(j\wedge(J-j))^{1/4},\qquad
 \rho_J=2-\frac{\log J}{J},
\]
where $z>0$ is fixed. Early crossings are controlled separately by a
coarser net of representatives. On the event with no crossing, the
Galton--Watson ballot estimate retains the cost of staying below the barrier,
rather than merely reaching a high endpoint.

To see where that cost enters, stop the traversal description at
$n=J-h$, with terminal radius $M=e^h$. Consider a cell in source shell
$\beta\log J+O(1)\le k\le C_0\log J$. The joint event that the prefix
stays below the barrier and $\sqrt{2T_n}\in[q,q+1]$ has probability at most
\begin{equation}\label{eq:ideas-ballot}
 C(k+1)^C(1+p)n^{-2}e^{-q^2/(2n)},\qquad
 p=B_n+1-q.
\end{equation}
This statement applies to the central range $\epsilon_0n\le q\le B_n+1$;
it is not a probability conditional on the prefix being good.
Extinction has substantial probability and is treated separately. The
factor $n^{-2}$ in \eqref{eq:ideas-ballot} follows by combining the
source law with the square-root prefactor in the ballot theorem. The
application-specific identity is
\[
 \frac{x^2}{m}+\frac{(q-x)^2}{n-m}
 =\frac{q^2}{n}+\frac{n(x-mq/n)^2}{m(n-m)},\qquad m=k+1.
\]
It combines the cost of the source height $x$ with the cost of its
descendants without losing an exponential factor.

\subsection{Keeping the terminal local times on the correct conditional law}

The traversal endpoint does not determine the local times inside the
terminal cell. We erase each completed outward path from radius $M$ to
radius $eM$, retaining its entrance and directed exit edge. Given the
exterior path and these endpoints, the erased paths are independent killed
bridges. If $b$ is their number and $X_i(x)$ their occupation at $x$, the
remaining quantity is $\max_x\sum_{i\le b}X_i(x)$.

Conditioning on a complete radial word also constrains the future of the
endpoint sequence. Consequently, a one-step harmonic-measure calculation
cannot simply be declared invariant under that conditioning. The proof
factors the entire word likelihood, including the last suffix and the
connectors between source intervals. It compares the true law with a law
obtained by normalizing the successive kernels in chronological order.
Uniform contraction controls dependence within each centered interval;
a fixed number of intervening connector kernels is allowed. This gives
both concentration of bridge moments and the exponential estimate needed
for the upper bound.

The central exact identity is a rank-one Green-function formula. For one
bridge with endpoint pair $\omega$, set $m_x(\omega)=E^\omega X(x)$ and
write $G$ for the Green function of its killing disk. Then
\[
 E^\omega e^{\theta X(x)}=1+r_\theta(x)m_x(\omega),\qquad
 r_\theta(x)=\frac{e^\theta-1}
 {1-(G(x,x)-1)(e^\theta-1)}.
\]
Averaging over the full endpoint law gives, to a controlled error,
$\log E[\exp\{\theta\sum_iX_i(x)\}\mid w]=bg r_\theta$.
For the upper bound choose $h=(C_h\log J)^4$ and initially
$r_0=(gn)^{-1}$. Since $G(x,x)=gh+O(1)$, the corresponding tilt satisfies
$\theta(r_0)u_{J,c}=2J-c\log J+O(1)$. A union bound over the $O(e^{2h})$
sites in one cell gives
\[
 P\left(\max_x\sum_iX_i(x)\ge u_{J,c}\mid w\right)
 \le Ce^{-2n}J^c e^{b/n}.
\]
Because $b=q^2/2+O(n)$, the factor $e^{b/n}$ cancels the exponential
in \eqref{eq:ideas-ballot}. This cancellation explains why the traversal
endpoint and the terminal maximum must be estimated together.

The displayed bound alone is insufficient: summing all endpoint bands
would lose powers of $J$. Varying the tilt around $r_0$ supplies the
additional factor
\[
 \exp\left\{-\kappa\min\left(
 \frac{(p-p_*)^2}{h},\frac{J^2}{h}\right)\right\},
 \qquad p_*=O(\log J).
\]
Its central sum, including the factor $1+p$, is
$O((1+|p_*|)\sqrt h+h)=J^{o(1)}$; the saturated part is negligible.
Thus one original cell contributes at most
$(k+1)^Ce^{-2n}J^{c-2+o(1)}$. There are $O(e^{2(n-k)})$ such cells in
shell $k$. Summing from $k=\beta\log J+O(1)$ produces exactly
$J^{-2\beta+c-2+o(1)}$. The remaining shells $k>C_0\log J$ contain
few enough sites that the elementary hitting and geometric-tail bounds
suffice, after choosing $C_0$ large.

\subsection{Producing a sufficiently large global maximum}

For the lower bound, the traversal construction must produce a cell with
\begin{equation}\label{eq:ideas-coarse-count}
 b=2L^2-2L\log L+O(L),\qquad L=J-h.
\end{equation}
Rosen's barrier construction provides such a root with positive probability
in the spherical Brownian model. Its angular restrictions make the
two-point sum finite. The passage to the lattice has to preserve this
existence event: an additive coupling error cannot be summed over
exponentially many rare candidate roots. Accordingly, the strong coupling
is applied to the union-existence event in one direction. The Brownian
pair estimate is used before the transfer. The lattice proof then uses a
first-moment bound and a separate conditional estimate for each root.

The last scales supply a contribution that cannot be discarded.
Conditionally on regular endpoint moments, the normalized sum of centered
bridge occupations has covariance $2gG(x,x')+O(1)$. A Gaussian
approximation and the maximum estimate for the discrete Gaussian free field
give
\begin{equation}\label{eq:ideas-terminal-lower}
 \max_x\sum_{i\le b}X_i(x)
 =gb+2\sqrt2\,g\sqrt b\,h
 +O_P\left(J\sqrt h+\frac{J\log J}{\log\log J}\right).
\end{equation}
Inserting \eqref{eq:ideas-coarse-count} shows the role of the second term:
\begin{align*}
 gb&=2gJ^2-4gJh-2gJ\log J
       +O(h^2+h\log J+J),\\
 2\sqrt2\,g\sqrt b\,h&=4gJh+O(h^2+h\log J).
\end{align*}
The terminal fluctuation restores the $4gJh$ lost by stopping the
traversal description early. The result is
$2gJ^2-2gJ\log J+o_P(J\log J)$.

Here the terminal scale is $h=(\log J)^p$ with $1<p<2$.
The inequality $p>1$ makes $M=e^h$ large enough for the relative
discretization errors. The inequality $p<2$ gives
$J\sqrt h=o(J\log J)$ and makes simultaneous moment concentration over
$O(M^4)$ observables possible. This is a different choice from the upper
bound, where no simultaneous screen of all these moments is needed.

Positive-probability disk success is enough. Partition a time interval of
length $e^{2j}$ into $j^a$ independent segments. Their effective logarithmic
radius is $j-(a/2)\log j+O(1)$, which costs $2a$ in the coefficient of
$gj\log j$. Choose $a>0$ within the available margin. Failure in every
segment has probability at most $e^{-cj^a}$, a summable bound. This yields
the deterministic-time lower estimate used for escape.

\subsection{Localizing the entire favorite set at the critical scale}

For critical recurrence, work at an independent geometric time $K$ with
$P(K=n)=(1-z)z^n$ and put $N=(1-z)^{-1}$. Choose a favorite by independent
continuous site priorities and reverse the path up to its last visit.
The reversed path consists of completed excursions from the selected site
and one terminal arm. Once the excursions, the forward suffix and the
priorities are fixed, being the selected favorite is exactly a collection
of upper bounds on the terminal arm's occupation at the other sites.
The arm is therefore a walk stopped when it returns to its root or exceeds
one of these capacities.

Its endpoint law is a positive mixture of stopped occupation measures.
Such a measure $\mu$, of mass $\beta$, satisfies
\[
 (I-zP)\mu=\alpha\delta_0-\nu,\qquad \nu\ge0,\qquad
 \alpha-\nu(\mathbb Z^2)=(1-z)\beta.
\]
Testing against a resolvent of the autocorrelation of a square gives
\begin{equation}\label{eq:ideas-resolvent-gain}
 \mu(B(0,C\sqrt N/f))
 \ge c\beta\frac{1+\log f}{f^2}.
\end{equation}
The Fourier transform of that test function is nonnegative, so its
resolvent is maximal at the origin. This makes the unknown nonnegative
exit measure $\nu$ usable. The two-dimensional integral is
$R^2(1+\log(N/R^2))$, which gives the logarithmic gain in
\eqref{eq:ideas-resolvent-gain} when $R\asymp\sqrt N/f$.
A restriction depending only on the exposed pieces makes the radius
proportional to $\sqrt K/f$ rather than $\sqrt N/f$.

Two remaining obstacles concern maximizers themselves. First, the selected
favorite can have distant ties. Insert the fixed loop
$(x,x+e_1,x,x+e_2,x)$ after its last visit. It raises the local time at $x$
by two and at each affected neighbor by one, making $x$ unique. The loop
is recovered from the last visit to the new unique favorite, so this is
an injection with a fixed positive probability cost. Second, a fresh
segment does not erase the old local-time profile $H$. If its own profile
$\ell$ has unique maximizer $y$, then
\[
 \operatorname*{arg\,max}(H+\ell)
 \subset\operatorname{supp}H\cup\{y\}.
\]
Thus it suffices to make the old range small on the new spatial scale;
the old local-time heights need not be bounded.

Choose block starting times $s_k=e^{J_k}$, with
$J_{k+1}-J_k=10\log J_k$, and take $f_k\asymp\sqrt{J_{k+1}}$.
The probability of a localized fresh singleton is at least
\[
 d_k\asymp\frac{\log J_k}{J_k},\qquad J_k\sim10k\log k.
\]
Hence $\sum_kd_k=\infty$. The logarithmic factor improves the available
success bound to order $1/k$; the weaker bound of order $1/J_k$ would
also have divergent sum for this particular schedule. The role of the
block choice is to make the old range negligible while preserving
divergence. More precisely, take
$N_k=e^{J_{k+1}}/J_{k+1}^2$ and retain clocks
$K_k\ge(d_k/8)N_k$. The old range is eventually at most
$\sqrt{s_k}J_k$, while $s_{k+1}/s_k\asymp J_k^{10}$ gives
\[
 \frac{\sqrt{s_k}J_k}{\sqrt{K_k/J_{k+1}}}
 \le\frac{C}{J_k^2\sqrt{\log J_k}}\longrightarrow0.
\]
At a smaller scale with $f_k\asymp J_k^\gamma$ and
$\gamma>1/2$, the probabilities furnished by the same estimate would
instead have summable lower bounds. This observation describes the
reach of this construction; the separate upper estimate proves escape.

Finally, the escape estimate with $c=2+2\delta$ gives a block probability
$j^{-2\gamma+2\delta+o(1)}$. For $\gamma>1/2$ choose $\delta>0$ so that
$2\gamma-2\delta>1$ and apply Borel--Cantelli. At $\gamma=1/2$, the
independent recurrence construction gives arbitrarily small normalized
favorite-set radii. These are the two summability calculations that place
the boundary at $1/2$.
